\documentclass[a4paper,11pt]{article}
\usepackage[utf8]{inputenc}
\usepackage{geometry}
\usepackage{xcolor}
\usepackage{listings}
\usepackage{hyperref}
\usepackage{graphicx}
\usepackage{amsmath}
\usepackage{amssymb}
\usepackage{tcolorbox}
\usepackage{float}

% Geometria della pagina
\geometry{top=2.5cm, bottom=2.5cm, left=2.5cm, right=2.5cm}

% Colori custom
\definecolor{ibmblue}{RGB}{0, 45, 156}
\definecolor{ibmred}{RGB}{250, 75, 76}
\definecolor{codegreen}{rgb}{0,0.6,0}
\definecolor{codegray}{rgb}{0.5,0.5,0.5}
\definecolor{codepurple}{rgb}{0.58,0,0.82}
\definecolor{backcolour}{rgb}{0.95,0.95,0.92}

% Setup per il codice Python
\lstdefinestyle{mystyle}{
    backgroundcolor=\color{backcolour},   
    commentstyle=\color{codegreen},
    keywordstyle=\color{magenta},
    numberstyle=\tiny\color{codegray},
    stringstyle=\color{codepurple},
    basicstyle=\ttfamily\footnotesize,
    breakatwhitespace=false,         
    breaklines=true,                 
    captionpos=b,                    
    keepspaces=true,                 
    numbers=left,                    
    numbersep=5pt,                  
    showspaces=false,                
    showstringspaces=false,
    showtabs=false,                  
    tabsize=2
}
\lstset{style=mystyle}

% Box personalizzati
\newtcolorbox{studentbox}{colback=blue!5!white,colframe=blue!75!black,title=Student Activity}
\newtcolorbox{theorybox}{colback=gray!5!white,colframe=gray!50!black,title=Theory Recap}
\newtcolorbox{warningbox}{colback=yellow!10!white,colframe=orange!80!black,title=Important: Qiskit Bit Ordering}
\newtcolorbox{composerbox}{colback=purple!5!white,colframe=purple!75!black,title=IBM Composer View (Q-Sphere)}
\newtcolorbox{deepdivebox}{colback=teal!5!white,colframe=teal!75!black,title=Deep Dive: Why does this matter?}

\title{\textbf{Laboratory Guide 2: Entanglement \& Multi-Qubit Systems}\\
\large The "Spooky Action at a Distance"}
\author{Course: Quantum Computing Foundation}
\date{}

\begin{document}

\maketitle

\section{Introduction and Objectives}
In Lab 1, we manipulated individual qubits. In this laboratory, we explore systems of multiple qubits and the phenomenon of \textbf{Entanglement}. As stated in Lesson QC05, the real power of quantum computation appears when we start combining qubits together.

\vspace{0.3cm}
\noindent \textbf{Prerequisites:}
\begin{itemize}
    \item Lesson QC05: \textit{Multiple Qubits}
    \item Lab 1: \textit{Introduction to QISKIT}
\end{itemize}

\noindent \textbf{Learning Objectives:}
\begin{enumerate}
    \item Understand the Controlled-NOT (CNOT) gate logic.
    \item Create a Bell State (Maximally Entangled State) using the Composer.
    \item Visualize correlations between qubits on the Q-Sphere.
    \item Verify entanglement statistics using Qiskit code.
    \item Extend entanglement to 3 qubits (GHZ State) and discuss scalability.
\end{enumerate}

\begin{warningbox}
\textbf{Reading Order in Qiskit}\\
Qiskit uses "Little Endian" ordering. This means the rightmost bit is qubit 0 ($q_0$).
\[ \text{State } |10\rangle \implies q_1 = 1, \quad q_0 = 0 \]
Always pay attention to which qubit is the \textit{Control} and which is the \textit{Target}.
\end{warningbox}

\newpage

\section{Part A: Logic \& Entanglement (IBM Quantum Composer)}
\textit{Platform: IBM Quantum Composer}

\subsection{Exercise 1: The CNOT Gate (Conditional Logic)}
\textbf{Context:} Lesson \textbf{QC05} introduces the CNOT gate. It acts on two qubits: a \textit{Control} and a \textit{Target}. It flips the Target qubit ONLY IF the Control qubit is in state $|1\rangle$.

\begin{studentbox}
\textbf{Action:}
\begin{enumerate}
    \item Use an \texttt{X} gate on $q_0$ (Control) to set it to $|1\rangle$. Leave $q_1$ (Target) at $|0\rangle$.
    \item Drag a \texttt{CNOT} gate (symbol $\oplus$ connected to a dot). Make sure the \textbf{dot (control)} is on $q_0$ and the \textbf{plus (target)} is on $q_1$.
    \item Observe the Output State.
\end{enumerate}
\textbf{Expected Effect:}
The state becomes $|11\rangle$ (Probability 100\% on "11"). The target flipped because the control was 1.
\end{studentbox}

\subsection{Exercise 2: Creating the Bell State (Entanglement)}
\textbf{Context:} Entanglement is a state where two particles become inextricably linked. The most famous example is the \textbf{Bell State} $|\Phi^+\rangle = \frac{|00\rangle + |11\rangle}{\sqrt{2}}$.

\begin{studentbox}
\textbf{Action:}
\begin{enumerate}
    \item Clear the circuit.
    \item Apply an \texttt{H} gate to $q_0$. (State is now superposition of 0 and 1).
    \item Apply a \texttt{CNOT} gate (Control $q_0$, Target $q_1$).
    \item Check the \textbf{Probabilities}.
\end{enumerate}
\textbf{Expected Effect:}
You should see two bars:
\begin{itemize}
    \item $\approx 50\%$ Probability for $|00\rangle$
    \item $\approx 50\%$ Probability for $|11\rangle$
    \item \textbf{0\% Probability} for $|01\rangle$ or $|10\rangle$.
\end{itemize}
\end{studentbox}

\begin{composerbox}
\textbf{Q-Sphere Visualization}\\
Look at the Q-Sphere. You will see two dots connected or symmetric. Unlike two separate qubits in superposition (which would show 4 dots: 00, 01, 10, 11), an entangled state only shows the correlated outcomes (00 and 11).
\end{composerbox}

\begin{deepdivebox}
\textbf{Why is Entanglement so Important? }\\
Entanglement is not just a weird phenomenon; it is the fundamental resource for two game-changing applications.

\vspace{0.2cm}
\textbf{1. Computational Power (The "Speedup")}\\
As we will see in \textbf{Lesson QC14} (Algorithms), entanglement allows a quantum computer to process a vast state space simultaneously. Here a brief description to some of those algorithms, that extensively makes use of the entanglement effect:
\begin{itemize}
    \item \textbf{Shor's Algorithm:} Uses entanglement to crack RSA encryption exponentially faster than classical computers.
    \item \textbf{Grover's Algorithm:} Uses entanglement to search through unsorted databases with quadratic speedup.
    \item \textit{Hardware Context:} Chips like \textbf{Google Willow} (2025) use 70+ entangled qubits to perform tasks impossible for classical supercomputers.
\end{itemize}

\vspace{0.2cm}
\textbf{2. Unhackable Security (Quantum Internet)}\\
Entanglement is the basis of \textbf{Quantum Key Distribution (QKD)}.
\begin{itemize}
    \item If two parties share entangled photons, any attempt by a hacker to "eavesdrop" (measure the line) destroys the entanglement instantly.
    \item This collapse reveals the intruder immediately, guaranteeing physical security.
    \item \textit{Hardware Context:} Companies like \textbf{IonQ} and \textbf{Quantinuum} are building the nodes for these quantum networks.
\end{itemize}
\end{deepdivebox}

\newpage

\section{Part B: Coding Entanglement (Qiskit)}

\subsection{Exercise 3: The "Hello World" of Quantum Code}
We will now write the Python code to generate the Bell State. This is the standard test to verify a quantum computer works.

\begin{tcolorbox}[colback=green!5!white,colframe=green!75!black,title=Coding Task: Bell State]
\begin{lstlisting}[language=Python]
from qiskit import QuantumCircuit
from qiskit_aer import AerSimulator
from qiskit.visualization import plot_histogram

# 1. Initialize Circuit with 2 Qubits
qc = QuantumCircuit(2, 2)

# 2. H on q0, then CNOT(0->1)
qc.h(0)
qc.cx(0, 1)

# 3. Measure both
qc.measure([0, 1], [0, 1])

# 4. Execute
backend = AerSimulator()
job = backend.run(qc, shots=1000)
print(f"Bell Results: {job.result().get_counts()}")
\end{lstlisting}
\end{tcolorbox}

\subsection{Exercise 4: Scaling Up - The GHZ State (3 Qubits)}
\textbf{Context:} State-of-the-art quantum processors (e.g., \textbf{Google Willow}, \textbf{IonQ Tempo}) rely on creating entanglement across many qubits. The simplest example of \textit{multipartite entanglement} is the \textbf{GHZ State}:
\[ |GHZ\rangle = \frac{|000\rangle + |111\rangle}{\sqrt{2}} \]

\begin{tcolorbox}[colback=yellow!10!white,colframe=orange!80!black,title=Concept Preview: What is Decoherence?]
You will study this formally in \textbf{Lesson QC06}, but for now, think of \textbf{Decoherence} as "environmental noise".
\begin{itemize}
    \item Just as a spinning top eventually falls due to friction, quantum states lose their superposition due to interaction with the outside world (heat, radiation).
    \item \textbf{The Rule:} The more qubits you entangle, the faster the system decoheres. This is why a 3-qubit GHZ state is much more fragile than a 2-qubit Bell state.
\end{itemize}
\end{tcolorbox}

\begin{tcolorbox}[colback=green!5!white,colframe=green!75!black,title=Coding Task: The 3-Qubit Entangler]
\begin{lstlisting}[language=Python]
# 1. Initialize 3 qubits
qc3 = QuantumCircuit(3, 3)

# 2. Entangle q0 and q1 (Bell State)
qc3.h(0)
qc3.cx(0, 1)

# 3. Entangle q2 (Chain reaction: q1 controls q2)
qc3.cx(1, 2)

# 4. Measure
qc3.measure([0, 1, 2], [0, 1, 2])

# 5. Execute
job3 = AerSimulator().run(qc3, shots=1000)
print(f"GHZ Results: {job3.result().get_counts()}")
\end{lstlisting}
\textbf{Expected Analysis:} Ideal results are '000' and '111'. In a real system, you would see errors (e.g., '001') appearing more frequently here than in the previous exercise, exactly due to the decoherence effect described above.
\end{tcolorbox}

\section{Part C: Assessment (Quiz)}

\textbf{Q1. The Logic of CNOT}\\
If the control qubit is in state $|0\rangle$ and the target qubit is in state $|1\rangle$, what does the CNOT gate do?
\begin{itemize}
    \item[A)] Flips the target to $|0\rangle$.
    \item[B)] Flips the control to $|1\rangle$.
    \item[C)] Nothing. The state remains $|01\rangle$ (Control is 0, so Target is untouched).
\end{itemize}

\vspace{0.3cm}

\textbf{Q2. Instantaneous Correlation}\\
You have created the Bell State $\frac{|00\rangle + |11\rangle}{\sqrt{2}}$. You measure the first qubit and get the result \textbf{"1"}. What is the state of the second qubit immediately after?
\begin{itemize}
    \item[A)] Still 50\% "0" and 50\% "1".
    \item[B)] Definitely "1". The wavefunction has collapsed to $|11\rangle$.
    \item[C)] Definitely "0".
\end{itemize}

\vspace{0.3cm}

\textbf{Q3. The Fragility of Large Systems}\\
In Exercise 4, we mentioned that "Decoherence" acts like noise. Why is the 3-qubit GHZ state harder to maintain than a single qubit?
\begin{itemize}
    \item[A)] Because Qiskit cannot handle 3 qubits.
    \item[B)] Because having more entangled particles increases the interaction surface with the environment, making the system lose information faster.
    \item[C)] Because 3 is an odd number.
\end{itemize}

\vspace{1cm}
\hrule
\vspace{0.2cm}
\footnotesize{\textbf{Solutions:} Q1: C; Q2: B; Q3: B (More qubits = Faster decoherence).}

\end{document}